\chapter{Modules and Test Definitions}\label{ch:modules}

A schedule is never one file on its own. It is part of a \textbf{Media Release}: a set of Lua
files delivered together and described by a \textbf{Test Definition}, a JSON file with the
extension \texttt{.atedef}.

\section{The Test Definition}

These are the files the manual's examples live in:

\begingroup\footnotesize
\begin{verbatim}
{
  "format": 1,
  "mediaNumber": "MR-MANUAL-TW",
  "name": "Test writer manual examples",
  "version": "1.0.0",
  "files": [
    { "path": "first.tslua",  "sha256": "a260d2b923819ff6..." },
    { "path": "channel.tslua", "hidden": true, "sha256": "1cecd8b3566e4783..." },
    { "path": "lib/rails.lua", "sha256": "392e385ab81dfb81..." },
    { "path": "modules/limits.lua", "module": "board_limits", "sha256": "ea8a21a8a7bef16b..." }
  ]
}
\end{verbatim}
\endgroup

This is an abbreviation of the real \texttt{manual.atedef} beside the examples: only the fields
that matter for reading it are shown, and the checksums are truncated.

\begin{itemize}
  \item \texttt{format} is \texttt{1}. An unknown version makes the Test Definition broken.
  \item \texttt{mediaNumber} is the release's identity. \texttt{name} and \texttt{version}
    (semantic version) are what other releases \texttt{depends} on.
  \item \texttt{files} lists every Lua file, with \texttt{path} and \texttt{sha256} --- the
    SHA-256 of the file's raw bytes, lowercase hex.
  \item \texttt{hidden} (only on a \texttt{.tslua}) keeps a schedule off the Run screen while
    still allowing it to be a Sub-schedule.
  \item \texttt{module} (only on a \texttt{.lua}) names it for \texttt{require}. A plain
    relative path derives its name: \texttt{lib/rails.lua} is \texttt{lib.rails}. An absolute
    path, or one containing \texttt{..}, must state \texttt{module}.
  \item Paths may be relative to the \texttt{.atedef} or absolute; forward slashes work on both
    Windows and Linux.
\end{itemize}

A Test Definition is broken if a listed file is missing, if a file has an extension that is
neither \texttt{.tslua} nor \texttt{.lua}, if two entries share a path or a module name, or if a
dependency cannot be resolved. A schedule that raises while loading is listed as broken with
that error; the rest of its release is unaffected.

Run \texttt{ate definition update <file.atedef>} to recompute every checksum in place.

\section{Checksums that do not match}

A mismatch is not a hard failure. ATE logs a WARNING naming the file and the expected and
computed checksums, loads the file anyway, and records both checksums with the Test Run.
Whether the differing bytes are stored is the \texttt{database.storeMismatchedFiles} config
setting: \texttt{ask} shows a confirmation at run start, \texttt{always} stores them,
\texttt{never} does not.

\section{\texttt{require}}

ATE replaces \texttt{require}. A module name resolves against the files in the release, then
against its dependencies, depth first in declaration order --- never against
\texttt{package.path}, which ATE empties. Asking for a module that is not listed is an error.
A module runs once per Test Run and its return value is cached.

The \texttt{ate.} prefix is reserved: \texttt{require("ate.devices")} is always the Device
params builder bundled with ATE, and a Test Definition may not list files under \texttt{ate.}.

\noindent\textbf{\texttt{modules.tslua}}\par
\lstinputlisting[language=ATELua]{../examples/test-writer/modules.tslua}

\noindent\textbf{\texttt{rails.lua}}\par
\lstinputlisting[language=ATELua]{../examples/test-writer/lib/rails.lua}

\noindent\textbf{\texttt{limits.lua}}\par
\lstinputlisting[language=ATELua]{../examples/test-writer/modules/limits.lua}

\cando{put a schedule and its helper modules in a Media Release, name the modules the way
\texttt{require} will find them, and check the checksums.}
